\begin{sddsprog}{sddslockin}
  \item \textbf{description:}
  \verb|sddslockin| performs lock-in (synchronous) detection of one or more columns of data as a function of a common
  independent variable, giving the amplitude and phase of each column at a list of specified frequencies.  For a
  frequency $f$ and a page with $N$ rows of data $(x_i, y_i)$, it computes
  \begin{equation}
    S = \frac{1}{N}\sum_i y_i \sin(2\pi f x_i), \qquad C = \frac{1}{N}\sum_i y_i \cos(2\pi f x_i),
  \end{equation}
  and reports the amplitude $2\sqrt{S^2+C^2}$ and the phase ${\rm atan2}(C,S)$, so that $y = A\sin(2\pi f x + \phi)$
  yields $A$ and $\phi$.  The frequency is in units of the inverse of the units of the independent column; for
  example, an independent column {\tt Pass} gives frequencies in cycles per pass (tunes).

  Each lock-in is a single term of a Fourier transform evaluated at exactly the requested frequency.  This makes the
  program the natural choice when one or a few frequencies are known in advance, for example the drive frequency of a
  frequency-response measurement: there is no need to select an FFT bin, the frequency need not fall on the FFT
  frequency grid, and the independent variable need not be uniformly spaced.  The estimate is exact when the data span
  an integer number of periods of the frequency with uniform sampling; otherwise it has leakage of order 1/(number of
  periods), and a nonzero mean of the data also leaks into the result unless \verb|-removeMean| is given.  The cost is
  proportional to the number of rows times the number of frequencies, so for a full spectrum \progref{sddsfft} is much
  faster; \verb|-sequence| and \verb|-equispaced| are best used to examine a limited frequency range in detail.

  Several dependent columns may be analyzed at once.  The sine and cosine are computed once per row and shared by all
  columns, so additional columns add little to the run time.
  \item \textbf{examples:}
  Find the amplitude and phase of the energy oscillation at the drive tune of a simulated frequency-response
  measurement, ignoring the first 4000 turns of transient:
\begin{verbatim}
sddsprocess run.w1 -pipe=out -filter=column,Pass,4000,1e12 \
  | sddslockin -pipe=in run.lockin -column=Pass,Cdelta \
      -atValues=9.342e-4 -removeMean
\end{verbatim}
  Analyze all centroid columns at two frequencies, using 8 threads:
\begin{verbatim}
sddslockin run.w1 run.lockin -column=Pass,C* \
  -atValues=9.342e-4,1.868e-3 -threads=8
\end{verbatim}
  Scan 1001 frequencies from 0 to 0.02 around a known line, at frequencies listed in another file, or from 0 to
  the Nyquist frequency in steps of $10^{-4}$:
\begin{verbatim}
sddslockin run.w1 run.scan -column=Pass,Cx,Cy -sequence=1001,0,0.02
sddslockin run.w1 run.scan -column=Pass,Cx -fileValues=tunes.sdds,nu
sddslockin run.w1 run.scan -column=Pass,Cx -equispaced=1e-4 -removeMean
\end{verbatim}
  \item \textbf{synopsis:}
\begin{verbatim}
sddslockin [-pipe[=input][,output]] [inputFile] [outputFile]
  -column=independentVariable,dependentVariable[,dependentVariable...]
  { -atValues=listOfFrequencies |
    -fileValues=filename,columnName |
    -sequence=points[,start,end] |
    -equispaced=spacing[,start,end] }
  [-removeMean] [-threads=number]
\end{verbatim}
  \item \textbf{files:}
  \emph{inputFile} is an SDDS file containing the independent column and the dependent columns.  If it contains
  multiple pages, each is analyzed separately and gives one page of \emph{outputFile}.

  \emph{outputFile} has one row per frequency, in a column {\tt Frequency}.  For each dependent column
  \emph{name}, it has columns \emph{name}{\tt Amplitude} (in the units of \emph{name}) and \emph{name}{\tt Phase}
  (in radians), and a parameter \emph{name}{\tt MeanRemoved} giving the mean subtracted with \verb|-removeMean| (zero
  otherwise).  The parameters {\tt LockinIndependentVariable} and {\tt LockinRows} give the name of the independent
  column and the number of rows analyzed.  All parameters of \emph{inputFile} are transferred to \emph{outputFile}.
  \item \textbf{switches:}
    \begin{itemize}
      \item \verb|-pipe[=input][,output]| --- The standard SDDS Toolkit pipe option.
      \item {\tt -column={\em independentVariable},{\em dependentVariable}[,{\em dependentVariable}...]} --- Specifies
      the independent column, followed by one or more dependent columns.  The dependent columns may be given as names,
      as wildcard patterns, or as any combination of the two.  Every numeric column that matches any entry is analyzed
      once, in the order given; a wildcard never selects the independent column or a non-numeric column.  A
      dependent column named explicitly must exist, but a wildcard pattern may match nothing provided that some
      column is selected.
      \item {\tt -atValues={\em listOfFrequencies}} --- Specifies a comma-separated list of frequencies.
      \item {\tt -fileValues={\em filename},{\em columnName}} --- Specifies that the frequencies are the values in
      column {\em columnName} of the first page of SDDS file {\em filename}.
      \item {\tt -sequence={\em points}[,{\em start},{\em end}]} --- Specifies {\em points} equally spaced
      frequencies, at least two.  If {\em start} and {\em end} are given, they specify the range, with {\em end}
      greater than {\em start}.  If they are not given, the range is from 0 to the Nyquist frequency of each page,
      $1/(2\Delta x)$, where $\Delta x$ is the mean spacing of the independent column; the frequencies may then differ
      from page to page.
      \item {\tt -equispaced={\em spacing}[,{\em start},{\em end}]} --- Specifies frequencies separated by
      {\em spacing}.  The range is determined as for \verb|-sequence|.
      \item \verb|-removeMean| --- Subtracts the mean of each dependent column, page by page, before detection, so
      that it cannot leak into the result.  This matters when the data do not span an integer number of periods.
      \item {\tt -threads={\em number}} --- Specifies the number of threads to use (default 1).  When there are at
      least as many frequencies as threads, the frequencies are divided among the threads; otherwise the sums over rows
      are.  With a single frequency the run time is usually dominated by reading the file, so threads help most with
      many frequencies.
    \end{itemize}
  \item \textbf{see also:}
    \begin{itemize}
      \item \progref{sddsfft}
      \item \progref{sddsnaff}
      \item \progref{sddssinefit}
    \end{itemize}
  \item \textbf{author:} Claude Code, M. Borland, ANL/APS.
\end{sddsprog}
